
\documentclass[twocolumn]{article}
\usepackage[top=1.1in, left=0.85in, right=0.85in]{geometry}

% \usepackage{supertabular}
\usepackage{amsthm}
% \usepackage{relsize}
\usepackage{amsmath}
\usepackage{amssymb}
% \usepackage{code}
\usepackage{graphicx}
\usepackage{fancyvrb}
\usepackage{url}
\usepackage{textcomp}

\pagestyle{empty}

\newcommand\comment[1]{}

\newcommand\st{$^{\mathrm{st}}$}
\newcommand\nd{$^{\mathrm{nd}}$}
\newcommand\rd{$^{\mathrm{rd}}$}
\renewcommand\th{$^{\mathrm{th}}$}
\newcommand\tm{$^{\mbox{\tiny \textsc{tm}}}$}

% nice fractions
\newcommand\sfrac[2]{{}\,$^{#1}$\!/{}\!$_{#2}$}

% \newcommand\citef[1]{\addtocounter{footnote}{1}\footnotetext{\cite{#1}}\ensuremath{^{\mbox{\footnotesize [\thefootnote]}}}}

\newcommand\dirham{\protect \includegraphics[width=1em]{dirham}}

% \usepackage{ulem}
% go back to italics for emphasis, though
% \normalem

\begin{document} 

% Academic Advancement Advice: Adopt A A Authorship
% Academic Advancement Advice: Author Articles as A A

\title{Academic Advancement Advice: Author Articles as A.~A.}
\author{A.~A. \\
  (tom7.org foundation)
  \thanks{
    Copyright \copyright\ 2018 the Regents of the Wikiplia Foundation.
    Appears in SIGBOVIK 2018 with the abecedarian advantage of the
    Association for Aomputational Aeresy; {\em IEEEEEE!} press,
    Verlag--Verlag volume no.~{\tt std::numeric\_limits<float>::lowest()}.
    1 \dirham
} }

\renewcommand\>{$>$}
\newcommand\<{$<$}

\date{1 April 2018}

\maketitle

\section*{Abstract}
proof at last
    1\ \dirham

\vspace{1em}
{\noindent \small {\bf Keywords}:
  A A A A A A
}

\section{Introduction}

Those with names falling in the dead-middle of the alphabet---like as
a random example, Murphy---have suspected since early life a
disadvantage. For example, when the teacher lines up the students for
pie, in order to be completely fair and avoid bias, orders them by the
perfectly objective criterion of alphabetical order. Occasionally, the
teacher realizes that the same students always get pie first, and
corrects this bias by using {\em reverse} alphabetical order.
Sometimes the teacher uses alphabetical or reverse-alphabetical by
first name, or sometimes last name. All variations of this idea are of
course unfair, as is clear to first- and second-graders, who are
powerless to rebel without (a) status in the K-12 hierarchy (b) the
statistical language in which to formalize their objection and (c)
access to a prestigious conference with permissive publication
standards.

This paper demonstrates that even within academic paper publishing
(which one could argue is obsessed with bias, prestigious conferences,
statistical language for formalizing objections, and status in the
K-12 hierarchy), this unfairness persists. Specifically, we find that
authors whose names come alphabetically earlier have higher citation
rates.

\section{The Data}

I obtained a FISA warrant based solely on a salacious and unverified
memo funded by SIGBOVIK's political opponents. The warrant was for the
AMiner Open Academic Graph, which consists of summary information for
154,771,162 papers in about 130 gigabytes of JSON data. Then I embarked
on one of the Classic Journeys of Computer Science.

\subsection{Relatable Computer Science Activities}

The task is apparently simple: Parse 130 gigabytes of JSON data to
extract the author names and citations from the papers and count them.
To begin my journey, I selected a C++ JSON library from 39
alternatives listed on a benchmark page.\cite{yip2018native} The first
criterion I ordered by was ``correctness,'' for which only two parsers
achieved the highest correctness score, $100\%$. I selected the slower
(i.e., more methodical) of the two parsers, since I did not want to do
shoddy work.

Next I spent a nice Saturday morning integrating this library into my
private build environment, making sure its open source license was
properly declared and so on, while I waited for the 130 gigabytes of
JSON to download. Miraculously, this was straightforward. I was able
to finish writing the code to process the articles before they even
finished downloading. However, upon unzipping, I found that each file
contained one million articles, not as a single JSON array but as a
series of adjacent objects. The JSON library allowed for reading JSON
from files, but {\em not} e.g. for reading a single record from a file pointer
repeatedly. So, I needed to do my own ad hoc parsing before feeding
each record to the JSON library as strings. OK. Now when I run it, it
prints
\begin{verbatim}
  [Segmentation fault
\end{verbatim}

This is not my first segmentation fault---let me tell you!---so I knew
what to do. I debugged the program, adding assertions using a
\verb+CHECK+ macro that aborts if its condition is violated. I
eventually determined that the bug must be inside the JSON library
itself, since the segmentation fault happened between two
\verb+CHECK+s where the only code being run was JSON library code. So,
deciding that perhaps its $100\%$ correctness benchmark score was
overstated, I downloaded and repeated the process with the other,
faster (i.e., more slapdash) parser. Of course this parser has a
different API so I needed to rewrite my program. (The program is very
simple so this was the fastest step; however, getting the library
properly incorporated and compiling and then learning its API was
still a computer science chore.) Once my program was rid of the old
crashy library, I ran it, and it produced this result
\begin{verbatim}
  [Segmentation fault
\end{verbatim}
Since these two libraries produced the same result, either
segmentation violation must be part of the JSON specification, or the
bug is instead elsewhere in the program. After reducing the program to
a fairly simple test case, I determined that it must be some kind of
``Heisenbug'' whereby an earlier piece of code was corrupting the heap
or similar, causing a later crash only when present. This was not my
first Heisenbug---let me tell you!---so I knew what to do. However,
not being able to use Valgrind on Windows (did I mention that I'm
using Cygwin and mingw64 on the eponymous Windows~7 by the way? Why am
I doing that?), I booted Linux in a virtual machine, installed the 87
security updates that had been released since I last booted this
virtual machine a few months ago, and compiled my program there. I
then needed to move one of the JSON files to the virtual computer to
test, without having to download 130 GB from the internet again, but
the file was apparently too large to transfer using the mysterious
``Drag and drop to virtual machine'' functionality. I recalled that
VirtualBox can mount a shared folder from the host machine if you edit
the \verb+/etc/fstab+ file to contain the correct arcane magic, and
that I had performed this ritual before, but then had commented-out
that line in ``rescue mode'' because I also learned that having this
line present during the boot process causes a kernel panic. So I
temporarily reinstated the line and mounted the drive and copied the
gigantic file into the virtual machine and then removed the dangerous
magic from the f-stab for f-safety.

Next I ran Valgrind on this minimal test case and it worked as
expected with no unclean behavior detected and no crashing.

Since Linux was no help debugging, I returned to Windows~7 and tried
single-step debugging. Here I discovered that the crash in my minimal
test case was occurring when using \verb+cerr+. Specifically, the
second thing sent to \verb+cerr+ would produce a segmentation fault.
I then recalled that the \verb+CHECK+ macro uses \verb+cerr+ and
also found code I had written in the Year 2014 

The original ``bug'' was a violated assumption: I assumed,
optimistically, that if the article had an {\tt author} field then
the author would be a string. I accounted for a missing author,
an author that was the empty string, but not {\tt author: null}.

multiple hypothesized causes:

 - When doing a 


300 GB of json woes
 

\subsection{Alliterative Articles}

Adopting an all-'A' affect also advances articles.
Awesome!

strange authors:

capinha para celular horizonte artificial iphone

\bibliographystyle{IEEEtran}
% \bibliographycomment{} % nothing
\bibliography{paper}

\end{document}
